\documentclass[a4paper,12pt,twoside,openright]{report}
\usepackage{graphicx}
\usepackage[hidelinks]{hyperref}
\usepackage[T1]{fontenc}
\usepackage[utf8]{inputenc}
\usepackage{setspace}
\usepackage{biblatex}
\usepackage[italian]{babel}
\usepackage{fancyhdr}
\usepackage{mathtools}

\preto{\section}{\filbreak}

\title{Applicazione per dispositivi mobili basata sul principio della bisimulazione}
\author{Lanzi Mattia}

\pagestyle{fancy}
\setlength{\headsep}{0.35in}
\let\MakeUppercase\relax
\renewcommand\sectionmark[1]{}
\setcounter{secnumdepth}{3}
\setcounter{tocdepth}{3}

\makeatletter
\let\inserttitle\@title
\let\insertauthor\@author
\makeatother

\begin{document}

\begin{titlepage}
\noindent
\begin{minipage}[t]{0.19\textwidth}
\vspace{-4mm}{\includegraphics[scale=1.15]{logo_unimib.pdf}}
\end{minipage}

\begin{minipage}[t]{0.81\textwidth}
	{
                {\textsc{Università degli Studi di Milano - Bicocca}} \\
                \textbf{Scuola di Scienze} \\
                \textbf{Dipartimento di Informatica, Sistemistica e Comunicazione} \\
                \textbf{Corso di laurea in Informatica} \\
                \par
        }
\end{minipage}
 
\vspace{40mm}
        
\begin{center}
	 {\LARGE{
                    \textbf{\inserttitle}
                    \par
            }}
        \end{center}
        
        \vspace{50mm}

        \noindent
        {\large \textbf{Relatore:} Prof.ssa Micucci Daniela} \\

        \noindent
        {\large \textbf{Primo Correlatore:} Prof.Ginelli Davide}

	\noindent
        	{\large \textbf{Secondo Correlatore:} Prof.Bernardinello Luca}
        
        \vspace{15mm}

        \begin{flushright}
            {\large \textbf{Relazione della prova finale di:}} \\
            \large{\insertauthor} \\
            \large{Matricola 827502} 
        \end{flushright}
        
        \vspace{40mm}
        \begin{center}
            {\large{\bf Anno Accademico 2020-2021}}
        \end{center}     
\end{titlepage}

\tableofcontents

\chapter*{Introduzione}
Il progetto di tesi consiste nello sviluppo di un'applicazione per dispositivi mobili che utilizzano una versione del sistema operativo Android.\\
L'applicazione dal titolo \textbf{Bisimulazione} permette di sperimentare un gioco basato sul principio della bisimulazione.\\
L'obiettivo della tesi è la realizzazione di una app Android che sia da supporto didattico agli studenti che affrontano lo studio della bisimulazione; in particolare, l'app sviluppata supporta lo svolgimento di un gioco dove attaccante e difensore si affrontano navigando un grafo fino a che sono disponibili delle mosse valide.\\
L'applicazione Bisimulazione ha dunque uno scopo prettamente didattico in quanto il suo fine è quello di essere utilizzato per comprendere il meccanismo della bisimulazione all'interno del corso xxx tenuto dal Prof.Bernardinello.

\chapter{Bisimulazione}
Un \textit{sistema a transizione di stati} (dall'inglese \textit{labelled transition system, LTS}), talvolta chiamato \textit{transition graph}, è una tripla (\textbf{Proc}, \textbf{Act}, \{ $\xrightarrow{\alpha}$ | $\alpha$ $\in$ \textbf{Act}\}) dove:

\begin{itemize}
\item \textbf{Proc} è un insieme di \textit{stati} (o \textit{processi});

\item \textbf{Act} è un insieme di \textit{azioni} (o \textit{etichette});

\item $\xrightarrow{\alpha}$ $\subseteq$ \textbf{Proc} x \textbf{Proc} è una \textit{relazione di transizione} per ogni $\alpha$ $\in$ \textbf{Act}.
\end{itemize}

Un sistema a transizioni di stati è \textit{finito} se i suoi insiemi di stati e azioni sono anch'essi finiti.\\
Assumiamo che \textit{P} e \textit{Q} siano stati in un LTS. Per ogni azione di $\alpha$, diciamo che \textit{P} $\xRightarrow{\alpha}$ \textit{Q} se una delle due condizioni è verificata:

\begin{itemize}
\item $\alpha$ $\neq$ $\tau$ e ci sono stati \textit{P'} e \textit{Q'} tali che\\
\centerline{\textit{P($\xrightarrow{\tau}$)*P'$\xrightarrow{\alpha}$Q'($\xrightarrow{\tau}$)*Q}}

\item oppure $\alpha$ = $\tau$ e P($\xrightarrow{\tau}$)*Q,
\end{itemize}

dove ($\xrightarrow{\tau}$)* indica la relazione riflessiva e transitiva $\xrightarrow{\tau}$.

\section{Dispositivi Android}
Secondo l'articolo riportato da Ericsson Mobility Report[4] sono 5.9 miliardi le persone nel mondo che possiedono un telefonino ed il numero di abbonamenti LTE risulta essere in crescita costante.\\
Tra questi emergono i dispositivi con sistema operativo Android, che risulta essere il sistema operativo più diffuso al mondo per quanto riguarda i dispositivi mobili (Figura 1.5). I fattori che hanno determinato questo primato di Android sono i prezzi contenuti per i mercati emergenti e la disponibilità di poter usufruire di migliaia di app per la salute e per l'istruzione.\\
Android, più di qualunque altra piattaforma per dispositivi mobili, aiuta le persone ad accedere alle informazioni e a più opportunità.[5]

\section{Obiettivo dello stage}
Come già anticipato, l'obiettivo dello stage è stato quello di implementare un'applicazione per dispositivi mobili che possa essere utilizzata come sostegno all'insegnamento e all'apprendimento del concetto della bisimulazione. Per lo sviluppo del progetto è stato scelto di programmare utilizzando l'\textbf{IDE} (\textit{integrated development environment})  \textbf{Android Studio}, un ambiente di sviluppo integrato adibito per la creazione di applicazioni Android, basato sul linguaggio \textbf{Java} per il \textit{backend} e sull'\textbf{XML} (\textit{extensible Markup Language}) per il \textit{frontend}. Android Studio offre inoltre un ampio supporto a serivizi di terze parti, in particolare ai serivizi di \textbf{Google}, tra cui \textbf{Firebase}, utilizzato principalmente come \textit{Realtime Database}.

\chapter{Organizzazione applicazione}
\section{Backend}
\subsection{Interazione con Firebase}
Firebase, servizio offerto da Google per l'utilizzo di funzionalità a supporto dello sviluppo mobile, interagisce con l'applicazione sotto tre aspetti principali:

\begin{enumerate}
\item registrazione all'applicazione;

\item autenticazione all'applicazione;

\item lettura e scrittura dati sul database.
\end{enumerate}

Nelle prossime sezioni verrà analizzato nel dettaglio ciascuno di questi tre aspetti.\\
Al fine di poter utilizzare il servizio di autenticazione e registrazione di Firebase è stato necessario dichiarare la dipendenza per la libreria Firebase Authentication nel modulo File Gradle a livello di app.\\
Allo stesso modo per poter usufruire delle funzionalità inerenti al database è stato necessario dichiarare la dipendenza per la libreria Firebase Database nel modulo File Gradle a livello di app.

\subsection{Registrazione}
La classe Java \textit{SignUp} si occupa della registrazione di un utente alla piattaforma Firebase che permetterà poi di autenticarsi al fine di poterne usufruire dei servizi.\\
La modalità scelta per la registrazione è la classica modalità che utilizza email e password. Uno degli attributi con modificatore privato di questa classe è \textit{auth}, un oggetto della classe FirebaseAuth; su questo oggetto, definito nel metodo \textit{onCreate()} (metodo che viene eseguito appena l'activity corrispondente alla classe viene lanciata), tramite il metodo \textit{getInstance()} che ritorna un istanza default della classe FirebaseApp, viene chiamato il metodo \textit{createUserWithEmailAndPassword()}.  Quest'ultima funzione legge come argomento due stringhe, l'indirizzo mail inserito dall'utente e la password, anch'essa inserita dall'utente, e prova a creare un nuovo account utente con l'indirzzo mail e la password fornita. Questo metodo lancia tre eccezzioni:

\begin{enumerate}
\item \textbf{FirebaseAuthWeakPasswordException};

\item \textbf{FirebaseAuthInvalidCredentialsException};

\item \textbf{FirebaseAuthUserCollisionException}.
\end{enumerate}

La prima eccezione viene lanciata nel momento in cui la password non è sufficientemente "forte", ovvero è inferiore ai 6 caratteri.\\
Per quanto riguarda l'eccezione \textit{FirebaseAuthInvalidCredentialsException} essa viene lanciata se l'indirizzo email inserito non segue il pattern corretto.\\
Infine l'ultima eccezione viene lanciata se è già presente un account con l'indirizzo mail inserito dall'utente.\\
Ogni eccezione viene gestita da un blocco catch specifico che provvede a mostrare all'utente messaggi differenti in modo da spiegare le cause degli errori.\\
 Se il processo va a buon fine, questa funzione autentica anche l'utente all'interno dell'applicazione; inoltre, le generalità inserite dall'utente vengono memorizzate all'interno del Realtime Database di Firebase: viene creato un oggetto \textit{map} di tipo HashMap<String, String> a cui vengono aggiunti tramite metodo put() tutte le informazioni necessarie e poi salvate nel databse. L'utente sarà aggiunto all'interno del nodo \textit{users}, come rappresentato nella figura XXX.

\subsection{Autenticazione}
L'autenticazione dell'utente all'interno dell'applicazione per utilizzare i servizi offerti da Firebase viene gestita dalla classe \textit{Login}. L'autenticazione avrà successo solo se l'utente in precedenza si è già registrato; in caso contrario non si potrà accedere e dunque non sarà reso possibile l'utilizzo dell'applicazione \textbf{Bisimulazione}.\\
A livello di codice l'autenticazione è gestita dal metodo \textit{signInWithEmailAndPassword()} che, così come per il metodo createUserWithEmailAndPassword() descritto nella sezione precedente, legge come argomento due stringhe, l'indirizzo mail e la password forniti dall'utente, e permette il login all'interno dell'applicazione. Questo metodo lancia due eccezioni:

\begin{itemize}
\item \textbf{FirebaseAuthInvalidUserException};

\item \textbf{FirebaseAuthInvalidCredentialsException}.
\end{itemize}

Le due eccezioni vengono lanciate rispettivamente se l'indirizzo mail dell'account utente non esiste oppure è stato disbilitato, o se la password inserita non è corretta.\\
Così come per le eccezioni descritte nella sezione precedente, anche in questo caso è presente un blocco catch per ognuna di esse.\\
La memorizzazione delle credenziali avviene tramite l'interfaccia \textit{SharedPreferences} che permette di accedere e modificare i dati a livello locale, non a livello di Realtime Database.
I dati collocati nelle SharedPreferences vengono salvati in un file XML contenuto nello storage interno, in particolare nella cartella \textit{shared\_prefs}. Il file in cui sono contenute le informazioni è chiamato "\textit{sharedPreferencesLogin}". In questo file vengono salvati i dati di accesso (indirizzo mail e password) dell'utente che saranno poi memorizzati per il successivo utilizzo.\\
Il salvataggio dei dati in questo file locale avviene solamente nel caso in cui l'utente abbia contrassegnato con un flag la checkbox accompagnata dalla descrizione "\textit{Ricordami}". In caso contrario, i dati non vengono salvati e se l'utente dovesse dimenticarsi della password utilizzata, è possibile effettuare un reset utilizzando sempre una funzionalità offferta da Firebase; nello specifico, per mezzo del metodo sendPasswordResetEmail() che prende come parametro l'indirizzo mail dell'utente è possibile effettuare questa operazione. Questo metodo attiva il backend di Firebase Authentication per inviare un'email di reimpostazione della password all'indirizzo email fornito che deve corrispondere ad un utente esistente già registrato a Bisimulazione.

\section{Realtime Database}
\textbf{Realtime Database} è un database ospitato nel cloud. I dati vengono archiviati come \textit{JSON} e sincronizzati in tempo reale con ogni client connesso, rimanendo disponibili fino a che l'app diventa offline. Con questo database è possibile archiviare e sincronizzare i dati con il database dal modello NoSQL.\\
Un aspetto fondamentale che ha favorito l'ultilizzo del Realtime Databsse come supporto all'applicazione è che, invece delle tipiche richieste HTTP, Firebase Realtime Database utilizza una sincronizzazione in tempo reale dei dati: ogni volta che i dati cambiano, qualsiasi dispositivo connesso riceve l'aggiornamento in pochi millisecondi. Questo aspetto è fondamentale per la gestione di una partita real time.\\

\subsection{Scrittura dati}
Per l'operazione di scrittura all'interno del database viene utlizzato il metodo \textit{setValue()} che permette di salvare il dato passato come argomento in un riferimento specifico, sostituendo qualsiasi dato esistente presente in quel percorso.\\
Questo metodo è stato utilizzato in maniera molto frequente e un esempio chiaro del suo funzionamento è possibile trovarlo nella classe Java \textbf{MatchmakingRoom}: qui quando viene creata una partita si passa al riferimento della partita stessa tutte le informazioni poi smistate nei vari nodi figlio, portando ad un risultato come quello in figura \textbf{xxx!!!!}.

\subsection{Lettura dati}
La lettura dei dati salvati all'interno del database è fondamentale per la creazione, lo sviluppo e il termine delle sfide in tempo reale.\\
Per leggere un valore presente nel database è stato utilizzato il meccanismo dei callback.\\
I dati memorizzati all'interno del Firebase Realtime Database sono recuperati collegando un ascoltatore asincrono (\textit{asynchronous listener}) che punta ad un riferimento del database. Il listener è attivato appena il dato viene trasmesso al database e in seguito si attiva ogni qual volta il dato cambia (si aggiorna o si elimina).\\
Innanzitutto viene dichiarata un'interfaccia, ovvero \textbf{DEFINIZIONE INTERFACCIA DA LIBRO DI TESTO}.
Di seguito la classe in cui si dovrà leggere il dato o i dati eredita l'interfaccia generata. A questo punto viene creato un metodo che legge come parametro il metodo callback definito nell'interfaccia; tramite un ascoltatore si legge il dato dal database che viene poi passato al callback. Una volta che il processo di lettura del dato termina, il callback rilascia il valore del dato letto nel metodo in cui è stato definito.\\
Per lo sviluppo di quest'applicazione mobile sono stati utilizzati esclusivamente due listener:

\begin{enumerate}
\item ascoltatore singolo (\textit{ListenerForSingleValueEvent});

\item ascoltatore persistente (\textit{ValueEventListener}).
\end{enumerate}

Un ascoltatore singolo permette di leggere i dati una volta sola; più nello specifico appena il callback viene chiamato e termina il suo processo, questo viene immediatamente rimosso.\\
Un esempio di utilizzo di un \textit{single value listener} è all'interno della classe Java \textbf{Table}: una volta letto il nome dell'attaccante e definito all'interno dell'info box non è più necessario mantenere un ascoltare su questo valore, quindi viene letto una volta sola e poi rimosso.\\
In maniera differente un ascoltatore persistente serve per tenere traccia di aggiornamenti ad uno specifico riferimento.\\
Un esempio di utilizzo di un \textit{value event listener} è possibile trovarlo sempre all'interno della classe Java \textbf{Table}: qui è necessario che venga letto il valore del riferimento \textit{turnOf} ogni volta che esso viene aggiornato per far sì che i giocatori che si sfidano sappiano in tempo reale a chi tocca eseguire una mossa.

\section{Frontend}
\subsection{Material Design}
typography

\subsection{Componenti Sign Up Activity}
Dal punto di vista dell'interfaccia grafica l'activity \textit{SignUp} presenta una serie di \textbf{TextView} ed \textbf{EditText} con lo scopo rispettivamente di informare l'utente sull'informazione richiesta che esso dovrà inserire e lo spazio necessario per eseguire tale operazione.

\subsection{Componenti Login Activity}
Questa activity, oltre a presentare le solite TextView ed EditText, mostra una \textit{checkbox} ovvero un componente che consente all'utente di selezionare un'opzione. In questo caso specifico  questa checkbox permette all'utente di memorizzare le proprie credenziali di accesso.\\
All'interno dell'activity è presente anche una \textit{bottom sheet}, ovvero un modello di interfaccia utente che viene utilizzato per visuallizzare una \textit{view} che si apre dal basso. Nell'applicazione ha funzione di reset della password in caso di dimenticanza.

\subsection{Componenti Home Fragment}
La schermata principale (\textit{Home}) costituisce l'unico \textit{Fragment} presente all'interno del progetto.\\
Qui sono presenti due \textit{cards} che forniscono all'utente informazioni sul proprio score personale e sulle partite attualmente attive.\\
Ciascuna card è una \textbf{MaterialCardView} che contiente al proprio interno TextView e Button. Un elemento presente all'interno di ciascuna card che poi ricorre anche in altri componenti all'interno delle varie interfacce è il \textbf{separator}, ovvero una semplice \textit{view} dalla larghezza di un solo 1dp.\\
In questo fragment troviamo anche un \textbf{FloatingActionButton}, ossia un pulsante circolare che permette di condividere la propria esperienza con l'applicazione sui social.

\subsection{Componenti Matchmaking Room Activity}
Questa activity è composta semplicemente da un button e un \textbf{ListView} posizionati all'interno di un \textit{Relative Layout}.\\
Una list view mostra un insieme di item che possono essere "scrollati" verticalmente e dove ogni item è posizionato immediatamente sotto l'item che lo precede nella lista.\\
Il componente ListView è stato implementato anche nell'activity \textit{ActivePlayers} per mostrare i giocatori attivi in quel momento.

\section{Canvas}
Al fine di disegnare i due grafi che compongono il tavolo da gioco dove si sfidano attaccante e difensore è stata utilizzata la classe \textbf{Canvas}.\\
Questa classe contiene le chiamate alle funzioni che permettono di disegnare.\\
Il metodo \textit{onDraw()} riceve come argomento un oggetto di tipo Canvas, che rappresenta la "tela" su cui disegnare gli oggetti, e al suo interno chiama due funzioni che hanno il compito di disegnare i nodi e gli archi dei singoli grafi.\\
Per quanto riguarda i nodi, questi vengono disegnati tramite il metodo \textit{drawCircle()} che legge quattro parametri:

\begin{enumerate}
\item la coordinata \textbf{x} del centro del nodo che sta per essere disegnato (di tipo \textit{float});

\item la coordinata \textbf{y} del centro del nodo che sta per essere disegnato (di tipo \textit{float});

\item il \textbf{raggio} del nodo che sta per essere disegnato (di tipo \textit{float});

\item un oggetto di tipo \textbf{Paint}.
\end{enumerate}

Quest'ultimo oggetto è il risultato della chiamata ad una funzione che riceve come parametro il colore del nodo e ritorna un oggetto con le informazioni sullo stile del nodo.\\ 
Gli archi invece sono disegnati tramite il metodo \textit{drawLine()} che riceve come parametri 5 argomenti:

\begin{enumerate}
\item la coordinata \textbf{x} del punto di partenza della linea;

\item la coordinata \textbf{y} del punto di partenza della linea;

\item la coordinata \textbf{x} del punto di arrivo della linea;

\item la coordinata \textbf{y} del punto di arrivo della linea;

\item un oggetto di tipo \textbf{Paint} esattamente come nel metodo \textit{drawCircle()}.
\end{enumerate}

Come punto di arrivo e punto di fine sono stati utilizzati come riferimento i nodi, facendo dei calcoli utilizzando il centro e il raggio.

\chapter{Partita in tempo reale}
\section{Lettura file di configurazione}
Al fine di poter creare nuove configurazioni di grafi da navigare durante la partita è stato utilizzato un meccanismo di lettura di un file di testo per generare i grafi.\\
Il file di testo, chiamato \textbf{config}, è posizionato nella cartella \textbf{assets}; tramite l'utilizzo della classe \textbf{InputStream} che permette di aprire il file di testo e della classe \textbf{BufferedReader} che da la possibilià di leggere il contenuto del file si è reso possibile creare nuove configurazioni senza dover modificare il codice sorgente.\\
All'interno del file sono descritte le specifiche da seguire per poter creare una nuova configurazione. Questo file deve essere modificato seguendo le specifiche e i pattern proposti; una eventuale modifica non corretta può portare a danneggiare le configurazioni già presenti.\\

\section{Interazione con Realtime Database}
\subsection{Preparazione}
Per preparare la sfida tra i due avversari viene inviata al database di Firebase una struttura di base che viene aggiornata durante la partita in base alle mosse degli sfidanti.\\
Una volta creata la partita dall'utente viene creato all'interno del nodo \textbf{rooms} una stanza che prende il nome del giocatore che l'ha creata, come in figura xxx.\\
I nodi figli di ogni singola stanza sono i seguenti:

\begin{itemize}
\item \textit{Player 1} che contiene a sua volta ulteriori nodi figli;

\item \textit{gameInProgress} inizializzato a \textbf{true};

\item \textit{lastMoveColour} inizializzato ad una stringa vuota;

\item \textit{leftGraph} che contiene nodi figli;

\item \textit{noMove} inizializzato a \textbf{false};

\item \textit{rightGraph} con lo stessa logica del nodo leftGraph;

\item \textit{show} inizializzato a \textbf{true};

\item \textit{specialColour};

\item \textit{turnOf} inizializzato alla stringa "attacker".
\end{itemize}

Il nodo padre, che indica il nome della stanza, acquisce un nodo figlio chiamato \textit{Player 2} nel momento in cui un avversario si unisce alla stanza. Questo nodo segue la stessa logica del nodo \textit{Player 1}. Questi nodi forniscono le informazioni sul nome dell'utente e il suo ruolo, attaccante o difensore.\\
Nel momento in cui il secondo giocatore si unisce alla partita, quindi nei momenti che ancora precedono l'inizio della stessa, il nodo \textit{show} assume valore \textbf{false}.
Il nodo \textit{show} se ha valore \textit{true} permette di mostrare la stanza della partita nella ListView, indicando che manca ancora un avversario. Quando questo nodo assume valore \textit{false} la stanza della partita non sarà più visibile all'interno della ListView.\\
Il nodo \textit{specialColour} assume il valore di una stringa corrispondente al colore di uno degli archi utilizzati all'interno del gioco (rosso, blu, verde, nero). Anch'esso non viene modificato durante lo svoglimento della partita.\\

\subsection{Svolgimento}
Durante la partita i nodi \textit{leftGraph} e \textit{rightGraph} vengono aggiornati costantemente.\\
Prendendo come esempio il nodo \textit{leftGraph} (il meccanismo per il grafo di destra è identico), esso contiene nodi figli che rappresentano tutti i nodi che compongono il grafo. A loro volta ogni nodo contiene due attributi: \textit{colour} e \textit{selected}. Questi due nodi vengono aggiornati ad ogni mossa dei giocatori, cambiando il loro valore se selezionati (colour assume valore blue e selected assume valore true).\\
Il nodo \textit{lastMoveColour} assume il valore del colore dell'arco percorso dal giocatore in seguito all'ultima mossa compiuta.\\
Il nodo \textit{turnOf} cambia ad ogni turno: inizialmente prende valore di "attacker" e ad ogni mossa eseguita, se valida, cambia il proprio valore passando a "defender" e viceversa.\\
Tutte queste informazioni sono contenute nella casella sottostante i tavoli di gioco.\\
In aggiunta è presente anche un attributo \textit{enabled} che, quando assume valore \textbf{false} rende il grafo inutilizzabile, come mostrato in figura xxx.\\
Infine il nodo \textit{gameInProgress} ha lo scopo di terminare la partita non appena esso assume valore \textbf{false}.

\chapter{Manuale utente}
\section{Manuale d'uso}
L'applicazione è costituita dalle seguenti activities:

\begin{itemize}
\item Login;

\item Registrazione;

\item Home / Main Activity;

\item Il mio profilo;

\item Gioca ora!;

\item FAQ;

\item Impostazioni;

\item Segnala un problema.
\end{itemize}

La funzione principale di questa applicazione è affrontare un avversario ad un gioco basato sul \textbf{principio della bisimulazione}. 

\section{Guida all'uso}
\subsection{Registrazione}
La schermata iniziale dell'applicazione \textbf{Bisimulazione} è rappresentata dalla schermata di \textbf{Login}. Questa schermata permette di autenticarsi e in 
seguito di utilizzare l'applicazione. L'autenticazione è possibile solo se già registrati. Se ancora non si è registrati, è necessario effettuare la registrazione.\\
Per poter effettuare la registrazione bisogna premere nel link a piè pagina contrassegnato dalla scritta \textit{Clicca qui}, immediatamente successivo
alla scritta \textit{Non sei ancora registrato?}.\\
A questo punto comparirà la schermata di \textbf{Registrazione}; qui basterà inserire i propri dati nelle caselle di testo dedicate e infine cliccare su \textit{Registrati}.
Una volta rilasciato il tasto \textit{Registrati}, i dati saranno inviati alla piattaforma Firebase e verrà aperta la schermata di \textbf{Login}.

\subsection{Login}
Di default quando si esegue l'applicazione viene mostrata la schermata di \textbf{Login}. Questa schermata ha come funzione principale quella di permettere agli utenti
di autenticarsi in maniera semplice e veloce al fine di poter usufruire dell'applicazione. Per autenticarsi è necessario inserire il proprio indirizzo mail utilizzato in 
fase di registrazione e la password scelta sempre durante la registrazione. A questo punto dovrà essere premuto il pulsante \textit{Login} e, una volta rilasciato,
si potranno verificare due casi:

\begin{enumerate}
\item l'autenticazione avviene \textbf{correttamente};

\item l'autenticazione \textbf{fallisce}.
\end{enumerate}
 
Nel primo caso viene mostrato un messaggio di avvenuta autenticazione e successivamente comparirà la schermata principale con titolo \textbf{Bisimulazione} (Main Activity).\\
Nel secondo caso viene stampato a video un messaggio di errore che non permette il passaggio alla schermata successiva, mantenendo l'utente sulla schermata
corrente. I motivi della mancata autenticazione possono essere molteplici a partire da un input scorretto fornito dall'utente.\\
Nel caso in cui si volessero memorizzare i propri dati di autenticazione è possibile spuntare con un flag la casella posizionata alla sinistra della scritta \textit{Ricordami}: in
questo modo ogni volta che verrà eseguita l'applicazione non bisognerà inserire di nuovo le proprie credenziali di accesso, bensì saranno già memorizzate e basterà
semplicemente cliccare sul pulsante \textit{Login}.\\
Inoltre se è stata dimenticata la propria password si può toccare il link descritto dal testo \textit{Clicca qui} immediatamente successivo alla scritta \textit{Password dimenticata?}.
A questo punto dal basso verrà mostrata una bottom sheet che invita l'utente ad inserire il proprio indirizzo mail utilizzato durante la registrazione; una volta digitato
 l'indirizzo mail e premuto il tasto \textit{Reset}, verrà inviata una mail all'indirizzo mail appena inserito con una nuova password da utilizzare al seguente accesso.\\

\subsection{Main Activity}
La schermata principale dal titolo \textbf{Bisimulazione} presenta due cards, tre pulsanti e il menù di navigazione.\\
Le due cards contengono informazioni riguardo l'utente, i giocatori e le partite in corso.\\ 
In particolare la prima card dal titolo \textit{Statistiche personali}  mostra le vittorie ottenute dall'utente sotto la scritta \textit{Numero di vittorie} e le sconfitte 
accumulate sotto la scritta \textit{Numero di sconfitte}. Inoltre questa card contiene un pulsante che presenta la dicitura \textit{Scopri di più} che se premuto apre la
schermata \textbf{Il mio profilo}.\\
La seconda card dal titolo \textit{Giocatori attivi} mostra il numero di giocatori attivi e le stanze da gioco attive al momento attuale.\\
Al centro sono presenti due pulsanti dal titolo: \textit{Gioca ora!!} e \textit{Giocatori attivi}. Il pulsante \textit{Gioca ora!} se cliccato porta alla pagina
\textbf{Stanza pre-partita}, mentre l'altro pulsante porta alla pagina \textbf{Giocatori attivi}.\\
In basso a destra è presente anche un floating action button che permette di condividere un messaggio con i propri contatti su differenti social networks: una volta cliccato
apparirà un messaggio che invita a condividere l'utilizzo dell'applicazione con i propri contatti. Se si clicca sul tasto \textit{Condividi}, si potrà proseguire condividendo
il messaggio, altrimenti se si tocca su \textit{Chiudi} la finestra verrà chiusa.\\
Infine questa schermata presenta lateralmente un menù di navigazione; per aprirlo basterà schiacciare sull'apposita icona e saranno mostrati le diverse activity
presenti all'interno dell'applicazione che possono essere raggiunte con un semplice tocco sull'item desiderato.

\subsection{Il mio profilo}
La schermata \textbf{Il mio profilo} presenta due sezioni: la prima sezione riguardanti le informazioni personali dell'utente dove è possibile trovare le generalità dell'utente
 che ha effettuato il login all'interno dell'applicazione quali nome, cognome e indirizzo mail utilizzato al momento della registrazione; nella seconda sezione intitolata \textit{Statistiche}
è mostrato il numero di vittorie e il numero di sconfitte accumulate dal giocatore.

\subsection{Giocatori attivi}
Questa activity mostra semplicemente una lista con tutti i giocatori attualmente attivi (ovvero che sono impegnati in una partita oppure sono in attesa di iniziarne una nuova).

\subsection{FAQ}
All'interno dell'applicazione è presente anche una pagina dedicata alle domande più frequenti (Frequently Asked Question) accessibile tramite menù di navigazione.
Questa pagina presenta alcune domande da cui si può ricavare la risposta semplicemente toccando sulla domanda a cui si è interessati. A questo punto si aprirà verso il basso
sotto la domanda un testo contenente la risposta alla domanda selezionata.

\subsection{Impostazioni}
La schermata \textbf{Impostazioni} permette esclusivamente di cambiare la lingua dell'applicazione da italiano a inglese o viceversa. Per eseguire questa operazione è necessario 
cliccare sulla scritta \textbf{Seleziona lingua...}; in seguito verrà mostrata la possibilità di scegliere tra \textit{Italiano} o \textit{Inglese}: per effettuare la propria scelta basterà toccare
sulla lingua desiderata e subito si tornerà alla schermata \textbf{Main Activity} con la nuova lingua selezionata.

\subsection{Segnala un problema}
Dal menù di navigazione è possibile raggiungere anche l'activity \textbf{Segnala un problema}. Qui è possibile segnalare eventuali problemi riscontrati durante l'utilizzo dell'applicazione
o fornire dei suggerimenti per migliorare l'applicazione; per fare ciò è necessario toccare all'interno della casella di testo dove è presente la scritta \textit{Inserisci qui il testo} e scrivere
il messaggio. Una volta terminato di scrivere il messaggio si deve cliccare sul tasto \textit{Invia} che, una volta rilasciato, mostrerà un messaggio di avvenuta spedizione del messaggio di testo.

\subsection{Stanza pre-partita}
L'activity \textbf{Stanza pre-partita} mostra tutte le partite create da altri giocatori che ancora stanno aspettando il proprio avversario. Per poter partecipare alla partita creata da un altro
giocatore si deve schiacciare su un elemento della lista con scritto \textit{Stanza di "Cognome"," Nome"}; a questo punto inizierà la partita in una nuova schermata contro l'avversario selezionato.\\
In questa activity è possibile anche creare una stanza propria cliccando sul pulsante \textit{Crea stanza}; appena viene cliccato questo pulsante si apre una nuova schermata con la partita.

\subsection{Gioco}
La partita tra i due avversari avviene nell'activity \textbf{Gioco}. Questa activity è composta da due tavoli: \textit{Tavolo sinistro} e \textit{Tavolo destro}.\\
All'interno di ogni tavolo è presente un \textbf{\textit{grafo}}. Il gioco consente di eseguire due azioni: toccare un nodo di uno dei due grafi oppure toccare il pulsante \textit{Stai fermo} che si trova
in fondo all'activity. Se viene toccato un nodo e la mossa eseguita dal giocatore è valida allora il nodo toccato si colorerà di blu, così come era colorato precedentemente il nodo selezionato.\\
Se il giocatore decide di non muoversi e questa mossa è consentita dalle regole del gioco, allora cambierà il turno senza che sui grafi avvenga alcun cambiamento.\\
In entrambi i casi, se la mossa eseguita rispetta le regole del gioco, viene aggiornato il turno di gioco e il colore dell'ultima mossa. In caso contrario, verrà mostrato un messaggio che suggerisce
al giocatore che ci sono ancora mosse disponibili; se questo non fosse vero, allora la partita terminerebbe mostrando un messaggio di \textit{Fine partita} che riporta l'utente alla \textbf{Main Activity}.
Sotto i due tavoli è presente una tabella che contiene informazioni utili riguardanti lo svolgimento del gioco; in particolare la casella fornisce informazioni su:

\begin{itemize}
\item \textbf\textit{attaccante};

\item \textbf\textit{difensore};

\item \textbf\textit{turno} del giocatore;

\item \textbf\textit{colore} dell'\textbf\textit{ultima mossa};

\item \textbf\textit{colore speciale}.
\end{itemize}

 Le regole del gioco possono essere trovate nel capitolo \textbf{3} "\textbf{Regole del gioco}".

\section{Regole}

Il gioco si svolge tra due giocatori, che chiameremo Attaccante (A) e Difensore (D).\\
Il terreno di gioco è formato da due grafi orientati, detti rispettivamente \textit{tavolo sinistro} e \textbf{tavolo destro}. Gli archi dei due grafi sono colorati di rosso, verde, blu e nero. All'inizio della
partita è proposto un colore che definiremo \textit{colore speciale}.\\
Una partita è composta da una sequenza di turni. Ad ogni turno muove l’attaccante che sceglie su quale tavolo compiere la mossa; e il difensore deve rispondere muovendo sul tavolo opposto.\\
Ci sono due tipi di mosse: mosse forti e mosse deboli. L’attaccante compie sempre mosse forti, il difensore mosse deboli. Una mossa forte consiste nello spostarsi dallo stato in cui si trova a uno stato raggiungibile percorrendo uno ed un solo arco del grafo. Diremo che il colore della mossa è il colore dell’arco percorso. Una mossa debole di colore \textit{x} consiste nello spostarsi lungo una sequenza di archi, uno dei quali di colore \textit{x}, mentre gli altri devono essere tutti del colore speciale. C'è un’eccezione: una mossa debole del colore speciale consiste in una sequenza di zero o più archi, tutti del colore speciale.\\
La partita può terminare in tre modi:

\begin{enumerate}
\item Dopo una mossa dell’attaccante, il difensore non può rispondere con una mossa dello stesso colore. L’attaccante vince.

\item Dopo una mossa del difensore, l’attaccante non ha mosse a disposizione su nessuno dei due tavoli. Il difensore vince.

\item Dopo una mossa del difensore, si torna a una configurazione già visitata nella stessa partita al termine di un turno. Il difensore vince.
\end{enumerate}

\chapter{Conclusioni}
In questa relazione di stage viene affrontato lo sviluppo di un'applicazione che mira ad essere di supporto per l'insegnamento didattico del principio della bisimulazione.\\
La soluzione offerta è dunque un gioco che permette di mettere in pratica le conoscenze acquisite a livello teorico su tale argomento contro un avversario che può essere più o meno competente in materia. Con questa soluzione si cerca di offrire una possibilità di testare le proprie competenze possibilmente divertendosi.\\
L'esperienza di stage interna all'università si è rivelata un importante occasione per accrescere le mie competenze in ambito di sviluppo Android, imparando ad utilizzare nuove librerie come le Canvas, aggiungendo al mio bagaglio personale nuovi meccanismi implementativi e perferzionando le abilità pregresse.

\chapter{Sviluppi futuri}
In questo capitolo vengono proposti dei possibili sviluppi futuri che possono rendere l'applicazione più completa e più fruibile.

\section{Scelta partita casuale}
Un primo aspetto che può rendere l'applicazione più facilmente utilizzabile è la possibilità di giocare una partita contro un avversario casuale. L'aggiunta di questa componente non rende più obbligatoria la scelta da parte dell'utente dell'avversario contro cui confrontarsi ma apre la possibilità a giocare una partita contro un avversario scelto in maniera randomica.

\section{Avversario CPU}
Attualmente è possibile affrontare esclusivamente un avversario reale in una partita online tramite l'interazione e le continue richieste con il database. In futuro si può sviluppare un'intelligenza artificiale da affrontare in mancanza di connessione Internet o semplicemente per testarsi in modalità singola.

\section{Animazioni}
Uno sviluppo futuro può riguardare la \textbf{UI} (\textit{user interface}) ovvero l'interfaccia grafica. Si possono aggiungere animazioni più intrattenenti quando un utente compie una mossa e dunque c'è un passaggio di stato all'interno del grafo scelto.

\chapter{Bibliografia}

\chapter{Sitografia}
\end{document}